% PCGS -- A&A research article. Figures: results/figures/nc (benchmark/figures/*.py); light JPEG copies in
% figs_light/ for review builds. Tables in tables/ (benchmark/paper_tables.py from the evaluation JSONs);
% Tables 1, 2, 4 and 7 are typed in (settings, image sets, volumes, development runs).
% Compile: tectonic manuscript.tex  (or pdflatex + bibtex)
\documentclass[fleqn,usenatbib]{aa}

\usepackage{graphicx}
\usepackage{txfonts}
\usepackage{amsmath}
\usepackage{booktabs}
\usepackage{placeins}
\usepackage[colorlinks=true,citecolor=blue,linkcolor=blue,urlcolor=blue]{hyperref}
% let figures and tables sit near their text: more floats per page, larger float areas
\renewcommand{\topfraction}{0.92}
\renewcommand{\dbltopfraction}{0.92}
\renewcommand{\bottomfraction}{0.6}
\renewcommand{\textfraction}{0.06}
\renewcommand{\floatpagefraction}{0.8}
\renewcommand{\dblfloatpagefraction}{0.8}
\setcounter{topnumber}{3}
\setcounter{dbltopnumber}{3}
\setcounter{totalnumber}{6}

\graphicspath{{../results/figures/nc/}}
\newif\iflightfigs \lightfigstrue   % JPEG copies in figs_light/ for a small review PDF
\newcommand{\figfile}[1]{\iflightfigs figs_light/#1.jpg\else #1.pdf\fi}

\newcommand{\pcgs}{PCGS}
\newcommand{\mzero}{\mu_0}
\newcommand{\vect}[1]{\boldsymbol{#1}}

\begin{document}

\title{PCGS: Photometrically constrained Gaussian splatting for sparse-view asteroid shape reconstruction}

\titlerunning{PCGS: Gaussian splatting for sparse-view asteroid shape reconstruction}
\authorrunning{X.-Y. Xiao \& F. He}

\author{Xin-Yu Xiao\inst{1,2}
        \and Fei He\inst{1,2}}

\institute{Institute of Geology and Geophysics, Chinese Academy of Sciences, Beijing 100029, China\\
           \email{xiaoxinyu24@mails.ucas.ac.cn}
           \and
           University of Chinese Academy of Sciences, Beijing 100049, China}

\date{Received ...; accepted ...}

\abstract
{Global shape models of small bodies are needed early in a rendezvous, when only a few dozen images are available.
Gaussian splatting can reconstruct surfaces from such image sets within tens of minutes, but it represents brightness
by view-dependent colours and ignores the solar illumination, which is known for every spacecraft image.}
{We aim to use the known illumination as a geometric constraint in Gaussian splatting and to evaluate the resulting
shape models against mission references, one of which is independent of images.}
{We present photometrically constrained Gaussian splatting (\pcgs). Each surfel carries a view-independent albedo,
which is shaded by a lunar-Lambert function under the Sun direction of each image, so that the observed brightness
constrains the surface orientation. At phase angles of 15$^\circ$ and above, cast shadows computed from the Sun's
viewpoint are added to the rendering; at lower phase angles, a total-variation prior is imposed on the albedo. A band
beyond the terminator is excluded from the silhouette loss to prevent the unlit part of the body from being removed,
and watertight meshes are fused from depth maps that are consistent across views. We applied \pcgs\ to five image sets
of Itokawa, Ryugu, and Bennu with phase angles of 7--60$^\circ$ and 40--50 training images each, and compared it with
four Gaussian splatting methods under identical inputs.}
{The mean distance between the \pcgs\ models and the mission shape models is 0.59~m for Itokawa and 0.62~m for Ryugu,
37\% and 15\% lower than that of the best competing method. At the registered scale, the volumes agree with the
references to within 0.3\%. Relative to the laser-altimeter model of Bennu, the mean distance is $0.62\pm0.06$~m,
compared with more than 3.7~m for the other methods; most of this difference results from the treatment of the night
side. The learned lunar-Lambert parameter decreases with increasing phase angle. On Itokawa and Ryugu, the method that
best reproduces the test images does not produce the most accurate shape.}
{Incorporating the known illumination into Gaussian splatting yields watertight global shape models with sub-metre
mean errors from about 50 images and about 30~min of training on a single GPU.}

\keywords{techniques: image processing --
          methods: numerical --
          planets and satellites: surfaces --
          minor planets, asteroids: individual: (25143) Itokawa --
          minor planets, asteroids: individual: (162173) Ryugu --
          minor planets, asteroids: individual: (101955) Bennu}

\maketitle

%--------------------------------------------------------------------
\section{Introduction}
\label{sec:intro}

The global shape of an asteroid is one of the first and most widely used products of a spacecraft encounter. Combined
with the mass, the volume yields the bulk density and hence the porosity, which identified both Itokawa
\citep{Fujiwara2006} and Ryugu \citep{Watanabe2019} as rubble piles. The shape also determines the surface gravity and
slopes, enters thermal and dynamical models, and supports navigation and the selection of landing and sampling sites
\citep{Lauretta2019}. Most mission shape models have been derived with two image-based techniques.
Stereophotogrammetry, including structure from motion and multi-view stereo (SfM-MVS), triangulates surface points that
are matched between images \citep{Preusker2015,Preusker2017,Watanabe2019}. Stereophotoclinometry (SPC) estimates the
slopes and albedo of small surface patches from their brightness in many images and links the patches through stereo
control points \citep{Gaskell2008,Jorda2016,Palmer2022}. Both techniques require hundreds to thousands of images with
repeated coverage, and SPC also requires considerable effort from analysts. In the first days of a rendezvous, however,
only a few dozen images from a narrow range of viewing geometries may be available, at the very time when a first shape
model is needed. Where a laser altimeter is available, it provides an independent reference; the SPC models of Bennu,
for example, were validated against the OSIRIS-REx Laser Altimeter \citep{Barnouin2019,AlAsad2021}, whose ranges also
yielded a global shape model \citep{Seabrook2022}.

Learned scene representations can reconstruct surfaces from such small image sets. Neural radiance fields
\citep{Mildenhall2020} and neural signed distance functions \citep{Wang2021NeuS} have been adapted to small bodies and
agree with mission shape models at the metre level \citep{ChenS2024,ChenH2024,ChenH2025}, but their training takes
hours to days. Gaussian splatting \citep{Kerbl2023} represents a scene by explicit Gaussian primitives that are
rasterised in real time, and its surface-oriented variants, such as two-dimensional Gaussian splatting
\citep[2DGS;][]{Huang2024}, Gaussian opacity fields \citep[GOF;][]{Yu2024GOF}, and PGSR \citep{Chen2024PGSR}, extract
meshes after tens of minutes of training. Asteroid-GS \citep{Zhang2025AsteroidGS} adapted Gaussian splatting to
asteroids by pruning Gaussians with masks, predicting colour with small networks that take the illumination as input,
and adding geometric regularisation; it reached accuracies comparable to stereophotogrammetry on Itokawa and Ryugu.

These methods inherit a model of appearance that was designed for terrestrial scenes under unknown illumination: the
brightness of a surface element is represented by a colour that varies freely with the viewing direction, often
complemented by per-image appearance codes \citep{MartinBrualla2021} or neural networks. Spacecraft images of an airless
body are formed differently. A surface of nearly uniform reflectance is illuminated by the Sun from a direction that is
known for every image, and the dependence of the brightness on the local surface orientation is the signal that
photoclinometry inverts for slopes \citep{Horn1990,Gaskell2008}. A free appearance model can absorb this signal instead
of exploiting it. In Earth observation, the Sun direction has been included in neural scene representations to model
shadows and solar illumination \citep{Derksen2021,Mari2022,Mari2023}, and EOGS introduced Sun-view shadow mapping into
Gaussian splatting \citep{SavantAira2025}. \citet{Nolan2026} replaced the colour of Gaussian splatting by a physically
based reflectance model for local terrain on Vesta and Ceres and obtained the most accurate normals with the
Lommel-Seeliger function. None of these approaches, however, has been applied to closed global shape models of small
bodies, for which the illumination raises two specific difficulties. At low phase angle, the shading contrast is weak,
and a darker surface element can be explained either by a slope facing away from the Sun or by a lower albedo. At high
phase angle, cast shadows become frequent, and a growing part of the body beyond the terminator is dark in every image;
a mask derived from such an image no longer outlines the body, and a silhouette constraint based on it removes the
unlit part of the body.

In this paper, we present photometrically constrained Gaussian splatting (\pcgs), which incorporates the known
illumination into 2DGS and addresses both difficulties. First, each surfel carries a view-independent albedo that is
shaded by the lunar-Lambert function \citep{McEwen1991} under the Sun direction of each image, so that the observed
brightness constrains the surface orientation. Second, the rendering depends on the phase angle: cast shadows rendered
from the Sun's viewpoint are added at moderate and high phase angles, and a total-variation prior on the albedo reduces
the ambiguity between albedo and slope at low phase angle. Third, a band beyond the terminator is excluded from the
silhouette loss, so that the unlit part of the body is not treated as empty space. We evaluate \pcgs\ on five image sets
of three asteroids with phase angles between 7$^\circ$ and 60$^\circ$. All compared methods receive identical inputs,
each image set is split three times at random into training and test images, and the reference model of Bennu was
derived from laser altimetry rather than from images.

%--------------------------------------------------------------------
\section{Method}
\label{sec:method}

\subsection{Overview}

Figure~\ref{fig:overview} illustrates the framework of \pcgs. The input consists of spacecraft images, camera poses
refined by bundle adjustment, a foreground mask for each image, and the Sun direction of each image in the body-fixed
frame. Because the masks are derived from the images, they contain only the illuminated part of the body. The surface
is represented by two-dimensional Gaussian surfels (Sect.~\ref{sec:2dgs}), which are rendered under the illumination of
each training image (Sect.~\ref{sec:render}) and optimised to reproduce it. The treatment of low phase angles and of the
night side is described in Sect.~\ref{sec:priors}, the training objective in Sect.~\ref{sec:losses}, and the extraction
of a watertight mesh in Sect.~\ref{sec:extraction}. \pcgs\ is built on the 2DGS framework and inherits its rasteriser,
its adaptive densification, and its normal-consistency regulariser.

\begin{figure*}
\centering
\includegraphics[width=\textwidth]{\figfile{fig1_overview}}
\caption{Overview of \pcgs. (a) The Ryugu reference model lit from the Sun direction of the held-out image in (c),
with the 60 cameras of the first split along their true directions (distances compressed; the spacecraft was about 20~km
from the asteroid): training views in blue (phase angle above 15$^\circ$, hence with cast shadows), held-out views in white, and the image of (c) in orange; yellow arrows indicate the Sun direction. (b) Spacecraft images with their SPICE Sun and camera
geometry train 2D Gaussian surfels that carry an albedo (inset: the trained surfels in a 120~m window, coloured by
their normal). Each image is rendered as $\hat I = A\,\Phi(\mzero,\mu)\,V$ (Eq.~\ref{eq:render}), and the surface
is extracted as a watertight mesh from depth that is consistent across views. Depending on the phase angle, the model adds either an albedo prior to the objective or cast shadows to the rendering. (c) A held-out Ryugu image (Hayabusa2 ONC-T, phase angle 19$^\circ$) decomposed by the trained
model into the rendered image, albedo, shading, normals, and residual. Renders of held-out views contain no cast
shadows, so the residual $\hat I-I$ is positive along the shadowed limb.}
\label{fig:overview}
\end{figure*}

\subsection{Gaussian surfels}
\label{sec:2dgs}

A 2DGS scene consists of planar Gaussian discs, or surfels, each defined by a centre, two tangent axes with scales, an
opacity, and a colour. For each pixel, the surfels intersected by the viewing ray are sorted by depth and
alpha-composited with the weights
\begin{equation}
w_k = o_k \prod_{j<k}\left(1-o_j\right),
\label{eq:weights}
\end{equation}
where $o_k$ is the opacity of surfel $k$ at the intersection point. The same weights give the rendered normal $\vect{n}$
and the accumulated opacity $O=\sum_k w_k$. We use the median depth, at which the accumulated opacity reaches one-half,
as the rendered depth; it is well suited to objects with sharp boundaries, such as an asteroid against the sky
\citep{Huang2024}.

\subsection{Photometric rendering}
\label{sec:render}

In \pcgs, the colour of surfel $k$ is a single albedo $a_k$ that does not depend on the viewing direction. The rendered
albedo $A=\sum_k w_k a_k$ is shaded per pixel as
\begin{equation}
\hat I = A\,\Phi(\mzero,\mu)\,V,
\label{eq:render}
\end{equation}
with the lunar-Lambert function \citep{McEwen1991}
\begin{equation}
\Phi(\mzero,\mu) = (1-L)\,\mzero + \frac{2L\,\mzero}{\mzero+\mu}.
\label{eq:ll}
\end{equation}
Here, $\mzero=\max(\vect{n}\cdot\vect{s},0)$ and $\mu=\max(\vect{n}\cdot\vect{v},10^{-3})$ are the cosines of the
incidence and emission angles, $\vect{s}$ is the unit vector towards the Sun computed from the SPICE kernels,
$\vect{v}$ is the unit vector towards the camera along the pixel ray, and $V\in[0,1]$ is the Sun visibility defined
below. The function varies between a Lambert surface ($L=0$) and the Lommel-Seeliger law ($L=1$) and equals unity at
normal incidence and emission. The partition parameter $L$ is learned as a single scalar per image set. Because the
shading is computed from the rendered normal, the photometric loss acts directly on the surface orientation. As the body
rotates under a nearly fixed Sun--spacecraft geometry, each surface element is observed under several Sun directions in
the body-fixed frame. With a single albedo shared by all views, the brightness of the element in these views constrains
its orientation, as in photometric stereo.

\paragraph{Cast shadows.} A local photometric function cannot explain shadows cast by boulders or by large-scale
topography. For images with a phase angle $\alpha\geq15^\circ$, we compute these shadows by shadow mapping
\citep{Williams1978}. The surfels are rendered with the same rasteriser from a virtual camera placed along $\vect{s}$ at
a distance of 50 body radii, which yields a Sun-view depth map $D_\odot$. For the surface point $\vect{x}_p$ seen in
pixel $p$, with depth $z_\odot(\vect{x}_p)$ from the Sun camera and projection $\pi_\odot(\vect{x}_p)$ into the Sun
view, the visibility is
\begin{equation}
V_p = \sigma\!\left(\frac{D_\odot\!\left(\pi_\odot(\vect{x}_p)\right) + b_p - z_\odot(\vect{x}_p)}{\tau}\right),
\label{eq:shadow}
\end{equation}
\begin{equation}
b_p = b_0 + \min\!\left(\Delta_\odot \tan \theta_p,\ 10\,\Delta_\odot\right),
\label{eq:bias}
\end{equation}
where $\sigma$ is the logistic function, $\tau$ controls the softness of the shadow edge, $\theta_p$ is the incidence
angle, and $\Delta_\odot$ is the pixel footprint of the Sun view. The slope-dependent bias $b_p$ prevents surfaces lit at
grazing incidence from shadowing themselves. Gradients reach both the shadowed point and the occluder, so that shadow
edges constrain the geometry. To prevent depth noise from producing spurious shadows, $V_p$ is applied only where the
observed image is darker than a fraction $\gamma$ of the unshadowed prediction $A\,\Phi$, and $V_p=1$ elsewhere; pixels
that are bright in the data are therefore never pulled into shadow. Below $15^\circ$, shadows are short and rare, and
the shadow map is not computed. Because this gating requires the observed image, renderings of test views contain no
cast shadows.

\subsection{Low phase angle and the night side}
\label{sec:priors}

\paragraph{Albedo prior.} At low phase angle the shading contrast is weak, and albedo and slope are difficult to
separate. To reduce this ambiguity, we penalise the total variation of the rendered albedo over the body in images with
$\alpha<15^\circ$:
\begin{equation}
\mathcal{L}_\mathrm{tv} = \frac{1}{|\Omega|}\sum_{p\in\Omega}\Big(\left|\tilde A_{p+\vect{e}_x}-\tilde A_p\right|
 + \left|\tilde A_{p+\vect{e}_y}-\tilde A_p\right|\Big),
\label{eq:tv}
\end{equation}
where $\tilde A=A/O$ is the opacity-normalised albedo and $\Omega$ is the set of body pixels with $O>0.5$. The prior
favours a piecewise-constant albedo and thus attributes smooth brightness variations to the geometry. Combined with the
cast shadows of Sect.~\ref{sec:render}, this defines a phase-dependent prior: the albedo prior is used below 15$^\circ$
and cast shadows from 15$^\circ$ upwards. The threshold was chosen during the development of the method
(Sect.~\ref{sec:rule}).

\paragraph{Night side.} The silhouette loss $\mathcal{L}_m$ is the binary cross-entropy between the accumulated opacity
$O$ and the foreground mask $M$. At high phase angle, a large part of the body beyond the terminator is dark in every
image and absent from the mask, and the silhouette loss would treat it as empty space. To retain this part of the body,
we exclude from the silhouette loss a band that extends the mask away from the projected Sun direction over a width
\begin{equation}
w_\mathrm{night} = 1.5\,r\,(1-\cos\alpha) + 5~\mathrm{pixels},
\label{eq:night}
\end{equation}
which is 1.5 times the maximum width of the unlit crescent of a sphere with the apparent radius $r$ of the body, plus a
margin. The first term amounts to 2\% of the apparent radius at $\alpha=10^\circ$ and to 35\% at $40^\circ$. The
photometric loss still applies within the band, so the night side remains constrained by the images in which it is lit.

\subsection{Multi-view consistency and training objective}
\label{sec:losses}

The shading constrains the orientation of the surface but not its absolute position. To constrain the position, we
follow PGSR \citep{Chen2024PGSR} and compare $7\times7$ patches of the observed images between each training view and a
neighbouring view, mapped by the homography of the rendered plane,
\begin{equation}
\vect{H} = \vect{K}_n\left(\vect{R}_{rn} - \frac{\vect{t}_{rn}\vect{n}^\top}{\delta}\right)\vect{K}_r^{-1},
\label{eq:homography}
\end{equation}
where $\vect{K}_r$ and $\vect{K}_n$ are the intrinsic matrices of the two views, $\vect{R}_{rn}$ and $\vect{t}_{rn}$
their relative rotation and translation, and $\vect{n}$ and $\delta$ the normal and the distance of the plane at the
patch centre. The loss $\mathcal{L}_\mathrm{ncc}$ is the mean of one minus the normalised cross-correlation over all
valid patches. Multi-view consistency thus constrains depth where the surface is textured, and shading constrains
orientation where it is not. The total training loss is
\begin{equation}
\mathcal{L} = (1-\lambda_\mathrm{s})\,\mathcal{L}_1 + \lambda_\mathrm{s}\,\mathcal{L}_\mathrm{D\text{-}SSIM}
 + \lambda_n\mathcal{L}_n + \lambda_m\mathcal{L}_m + \lambda_v\mathcal{L}_\mathrm{ncc} + \lambda_\mathrm{tv}\mathcal{L}_\mathrm{tv},
\label{eq:loss}
\end{equation}
where the first two terms compare $\hat I$ with the observed image through the absolute difference and the structural
dissimilarity \citep{Kerbl2023}, $\mathcal{L}_n$ is the normal-consistency term of 2DGS, and $\lambda_\mathrm{tv}=0$ for
images with $\alpha\geq15^\circ$. The weights are listed in Table~\ref{tab:settings}.

\subsection{Surface extraction}
\label{sec:extraction}

Surfels that float in front of the surface but still reproduce the images are a major source of error in meshes
extracted from Gaussian splatting. To suppress them, we keep only depth values on which several views agree. For each
training view, we render the median depth and the normal, and a pixel is kept if at least three of six neighbouring
views see its back-projected point within 1.5~m of their own rendered depth. The kept oriented points are meshed by
screened Poisson reconstruction \citep{Kazhdan2013}, and the largest connected component is retained. Poisson
reconstruction closes regions that no image constrains, such as the deep neck of Itokawa, but tends to pull them
inwards. We therefore re-solve, or fair, the vertices that lie more than 2~m from a truncated signed-distance fusion
\citep{Curless1996} of the same depths as a smooth surface of minimal thin-plate energy \citep{Botsch2004}, which keeps
the mesh watertight. Finally, the rendered median depth lies slightly in front of the true surface, by about 0.15 times
the pixel scale in our data. We remove this bias by offsetting the mesh along its normals by the median signed distance
of the SfM tie points.

\subsection{Implementation details}
\label{sec:implementation}

The surfels are initialised from the SfM tie points and optimised for 30\,000 iterations with the optimiser and the
densification schedule of 2DGS (Table~\ref{tab:settings}); the depth-distortion loss of 2DGS is not used. During the
first 3000 iterations, the rendered albedo is compared directly with the image to obtain a coarse geometry, after which
shading and cast shadows are switched on. The parameter $L$ is optimised through a logistic parametrisation with a
learning rate of $10^{-2}$, and the photometric loss is evaluated over the whole image with a black background. For the
multi-view term, 4096 pixels inside the foreground mask are sampled per iteration. The neighbouring view is drawn from
up to four views whose directions, seen from the body centre, differ by 2--25$^\circ$, and patches that face away from
this view or are occluded in it are discarded. The Sun camera is a perspective camera with a field of view of
3$^\circ$, which is close to orthographic, and no gradient is propagated through the shadow gating. The apparent radius
$r$ in Eq.~(\ref{eq:night}) is the 99th-percentile distance of the SfM tie points from their centroid, converted into
pixels with the camera distance. For the extraction, the six neighbouring views are those whose directions differ by
3--40$^\circ$, and every second pixel of the kept depth is used. The signed-distance fusion has a voxel size of 1~m and a
truncation distance of 4~m. The faired region extends two rings of vertices beyond the unsupported vertices, and all
other vertices are held fixed.

The public implementation of 2DGS places the principal point half a pixel away from the pixel convention of COLMAP
\citep{Schonberger2016}, which we use for the bundle adjustment, and its back-projection of depth adds a further half
pixel. We correct both offsets, so that renderings and cameras agree (Sect.~\ref{sec:protocol}). Training took
17--47~min per reconstruction, with a median of 27~min, on one NVIDIA A100 GPU of a shared computing node, and the mesh
extraction a further 4--10~min. The shadow map adds a second rasterisation per iteration.

\begin{table}
\caption{Settings of \pcgs, identical for all image sets.}
\label{tab:settings}
\centering
\footnotesize\setlength{\tabcolsep}{3pt}
\begin{tabular}{@{}ll@{}}
\toprule
Component & Setting \\
\midrule
Training & 30\,000 iterations; optimiser of 2DGS \\
Appearance & albedo (SH degree 0); $L$ learned, initial 0.9 \\
Shading, shadows & from iteration 3000 \\
Image loss & $\lambda_\mathrm{s}=0.2$ (D-SSIM) \\
Normal consistency & $\lambda_n=0.05$ from iteration 7000 \\
Silhouette & $\lambda_m=0.05$; night-side band, Eq.~(\ref{eq:night}) \\
Multi-view NCC & $\lambda_v=0.15$ from iteration 7000; 4096 \\
 & patches; $\leq$4 neighbours at 2--25$^\circ$ \\
Albedo TV & $\lambda_\mathrm{tv}=0.1$ for $\alpha<15^\circ$ \\
Cast shadows & $\alpha\geq15^\circ$; Sun view $1024^2$ pixels; \\
 & $\tau=0.5$~m, $b_0=3$~m, $\gamma=0.9$ \\
Extraction & $\geq$3 of 6 views agree within 1.5~m; \\
 & Poisson depth 11; fairing beyond 2~m \\
Scene extent & 1.1 $\times$ 99th-percentile SfM radius \\
\bottomrule
\end{tabular}
\end{table}
%--------------------------------------------------------------------
\section{Data and experimental setup}
\label{sec:setup}

\subsection{Image data}
\label{sec:images}

We used five image sets of three asteroids (Table~\ref{tab:data} and Fig.~\ref{fig:campaigns}). The Itokawa images were
acquired by the Asteroid Multi-band Imaging Camera \citep[AMICA;][]{Ishiguro2010} on Hayabusa at distances of
7.3--8.0~km \citep{Saito2006}. The Ryugu images were acquired by the telescopic Optical Navigation Camera
\citep[ONC-T;][]{Kameda2017} on Hayabusa2 during three hovering campaigns at about 20~km. The phase angle was nearly
constant within each campaign, at 7$^\circ$, 19$^\circ$, and 40$^\circ$, respectively, so the three image sets show the
same body under three illumination geometries. The Bennu images were acquired by MapCam \citep{Rizk2018} on OSIRIS-REx
during the Approach and Preliminary Survey phases, at distances of 9--38~km and phase angles of 26--60$^\circ$.

From each acquisition window, we selected 60 frames that show the whole body and are evenly spaced in sub-spacecraft
longitude; for the 40$^\circ$ Ryugu campaign, 48 frames were available. The Bennu frames were selected from the SPICE
geometry alone, so some of them show only part of the body. All images have $1024\times1024$ pixels. The calibrated
ONC-T and MapCam images were converted to 8 bits with one scale factor per image set, which preserves their relative
radiometry, and the AMICA images were used as raw 8-bit numbers. Each image set was split three times at random into 50
training and 10 test images, or 40 and 8 images at 40$^\circ$. \pcgs\ was trained on all three splits of Itokawa,
Ryugu at 19$^\circ$, and Bennu, on two splits at 40$^\circ$, and on one split at 7$^\circ$.

Four of the five image sets influenced the design of the method. The Itokawa and 19$^\circ$ Ryugu image sets were used
during development, the 40$^\circ$ Ryugu set informed the choice of the phase threshold, and a preliminary
reconstruction of Bennu without the night-side band motivated the band. Only the 7$^\circ$ Ryugu set played no part in
the development.

\begin{table*}
\caption{Image sets.}
\label{tab:data}
\centering
\small\setlength{\tabcolsep}{4pt}
\begin{tabular}{@{}llllcccc@{}}
\toprule
Body & Camera & Dates & Images & Distance (km) & Pixel scale (m) & Phase angle ($^\circ$) & Splits (training/held-out) \\
\midrule
Ryugu   & Hayabusa2 ONC-T  & 2019 Jan 15        & 60 & 20.18--20.26 & 2.18--2.19 & 7.1--7.3   & 3 $\times$ 50/10 \\
Itokawa & Hayabusa AMICA   & 2005 Sep 29--Oct 2 & 60 & 7.32--8.00   & 0.73--0.79 & 9.3--10.4  & 3 $\times$ 50/10 \\
Ryugu   & Hayabusa2 ONC-T  & 2018 Jul 10        & 60 & 19.90--20.09 & 2.15--2.17 & 18.5--18.7 & 3 $\times$ 50/10 \\
Ryugu   & Hayabusa2 ONC-T  & 2018 Aug 31        & 48 & 21.77--21.89 & 2.35--2.36 & 39.9--40.3 & 3 $\times$ 40/8 \\
Bennu   & OSIRIS-REx MapCam & 2018 Nov 30--Dec 13 & 60 & 9.07--38.31 & 0.62--2.60 & 26.1--59.9 & 3 $\times$ 50/10 \\
\bottomrule
\end{tabular}
\tablefoot{The pixel scale is the distance to the body centre divided by the focal length (10\,066.7, 9\,269.2,
and 14\,729.4 pixels for AMICA, ONC-T, and MapCam). Sub-spacecraft latitudes are within 5$^\circ$ of the equator
for Itokawa and Ryugu and between $-4^\circ$ and $13^\circ$ for Bennu; sub-solar latitudes lie between $-8^\circ$ and $+3^\circ$.}
\end{table*}

\begin{figure*}
\centering
\includegraphics[width=\textwidth]{\figfile{fig_campaigns}}
\caption{One held-out image of each image set, ordered by phase angle $\alpha$: (a) Ryugu, ONC-T, 2019 January 15; (b) Itokawa, AMICA, 2005 September 30; (c) Ryugu, 2018 July 10; (d) Ryugu, 2018 August 31; (e) Bennu, MapCam, 2018 December 13. Panels (a), (c), and (d) show the same body under three illumination geometries; at high phase angle (d, e) a large part of the body lies on the unlit night side.}
\label{fig:campaigns}
\end{figure*}

\subsection{Camera poses and foreground masks}
\label{sec:cameras}

Initial camera poses were computed from the SPICE kernels \citep{Acton1996}. We matched SIFT features between all image
pairs with COLMAP and refined the poses by bundle adjustment, starting from the SPICE poses with a reprojection-error
threshold that was reduced from 40 to 4 pixels. Incremental structure from motion without this initialisation failed for
these narrow-field cameras. A similarity transform maps the refined camera centres onto the SPICE centres, so the scene
is expressed in the body-fixed frame at the SPICE scale and the SPICE Sun vectors apply directly. The foreground masks
were obtained by thresholding each image, at 8~DN for AMICA and at 4\% of the 99.9th percentile of the brightness for
ONC-T and MapCam, keeping the largest connected region, and moving its edge to the half-maximum of the blurred limb. All
methods use the same poses and masks.

\subsection{Reference shape models and registration}
\label{sec:reference}

The Itokawa reference is the SPC model of \citet{Gaskell2008} with 3.1 million facets, and the Ryugu reference is the
SfM-MVS model of \citet{Watanabe2019} with 3 million facets. Both models were built from images, and the Itokawa model
shares images with our set. The Bennu reference is the global shape model v21 derived from the OSIRIS-REx Laser
Altimeter \citep[OLA;][]{Seabrook2022}, with 3.15 million facets and a mean edge length of 0.88~m; it is independent of
any image. For each image set, a single similarity transform from the camera frame to the reference frame is estimated
by aligning the SfM tie points to the reference surface with a trimmed iterative closest point (ICP) algorithm. Its scale
factor lies between 0.997 and 1.011 (Table~\ref{tab:volume}) and measures the agreement between the SPICE scale of the
camera frame and that of the reference. The same transform is applied to all reconstructions of an image set, and no
reconstruction is fitted to the reference.

\subsection{Compared methods}
\label{sec:baselines}

We compare \pcgs\ with four Gaussian splatting methods that extract surfaces: 2DGS, GOF, PGSR, and Asteroid-GS. Each
method is trained with its authors' implementation and meshed with the extraction they recommend; for 2DGS, we use its
configuration for object-centric scenes. All methods receive the same images, poses, masks, and splits. They also share
the same object extent, 1.1 times the 99th-percentile radius of the SfM tie points, which sets the spatial scale of
densification instead of the radius of the camera ring, which is about 30 times larger. Asteroid-GS is additionally
evaluated with the scene extent of its public code, a variant denoted Asteroid-GS$^\dagger$. GOF exceeded the available
GPU memory on Bennu. As a reference for image fidelity alone, we also trained Mip-Splatting \citep{Yu2024Mip}, a
novel-view synthesis method without surface constraints, on one split of Itokawa and fused a mesh from its rendered
depth. On the same computing node, training took 29--31~min per reconstruction for PGSR on Itokawa and Ryugu and
22--52~min for Asteroid-GS; all training times depend on the load of the node.

\subsection{Evaluation metrics}
\label{sec:metrics}

Following \citet{Zhang2025AsteroidGS}, we report the mean and the root mean square error (RMSE) of the distance from the
reconstructed vertices to the reference surface. Because these one-sided measures do not penalise missing surface, we
also compute symmetric measures on $N=10^6$ area-uniform samples $\mathcal{X}$ of the reconstruction and $\mathcal{Y}$
of the reference. With $d(\vect{x},\mathcal{Y})$ denoting the distance from a point to the other surface, the accuracy
and the completeness at a threshold $\epsilon$ are
\begin{equation}
P(\epsilon) = \frac{1}{N}\sum_{\vect{x}\in\mathcal{X}}\left[d(\vect{x},\mathcal{Y})<\epsilon\right],\qquad
R(\epsilon) = \frac{1}{N}\sum_{\vect{y}\in\mathcal{Y}}\left[d(\vect{y},\mathcal{X})<\epsilon\right],
\label{eq:pr}
\end{equation}
and the F-score is their harmonic mean; we use $\epsilon=1$~m. The Chamfer distance is the mean of the two directed mean
distances, and the normal error is the mean angle between the normal of each sample and the normal of its closest point
on the other surface. Image fidelity is measured on the test images by the peak signal-to-noise ratio (PSNR), the
structural similarity index (SSIM), and the learned perceptual image patch similarity (LPIPS)
\citep{Wang2004SSIM,Zhang2018LPIPS}, after fitting one gain and offset per image by least squares; this step is needed
because methods with per-image appearance codes have no code for an unseen image. We report the mean and the standard
deviation over the three splits. Because any two splits share at least 40 of their 50 training images, the standard
deviations describe the sensitivity to the image selection rather than the full uncertainty.

\subsection{Post-processing and registration}
\label{sec:protocol}

Every mesh, including those of \pcgs, receives the same reference-free post-processing: a translation and a normal
offset fitted to the SfM tie points, 20 iterations of Taubin smoothing \citep{Taubin1995}, and the selection of the
largest connected component. The alignment to the tie points matters because of a pixel-convention offset that affects
the Gaussian splatting codes. The projection of the 2DGS code, and of codes derived from the same rasteriser, assumes a
principal point at $(W-1)/2$ pixels for an image of width $W$, whereas COLMAP places the centre of pixel $i$ at $i+1/2$,
and the back-projection of depth uses yet another convention, $W/2$. Together, these conventions shift the rendering by
one pixel relative to the bundle-adjusted cameras. A synthetic test that back-projects a rendered depth map and
reprojects it with the COLMAP camera gives a residual of $(-1,-1)$~pixel before our correction and $10^{-3}$~pixel
after it. For PGSR and Asteroid-GS on Ryugu, the alignment to the tie points removes a translation of 2.15 and 1.9~m,
respectively, along the spin axis, consistently across splits. This corresponds to about one pixel along the image
vertical at the 2.2~m pixel scale, which is consistent with the convention offset, although we have not traced it within
each code. Without the alignment, the mean distances of these methods roughly double; for PGSR on Ryugu, the mean
distance is 1.51~m instead of 0.73~m. The alignment does not change the results of \pcgs. Table~\ref{tab:protocol}
compares the meshes as extracted, after the alignment, after smoothing, and after an ICP registration of each mesh to the
reference. \pcgs\ ranks first under every protocol, whereas the order of the baselines changes.

\begin{table*}[t]
\centering
\caption{Mean distance and F-score at 1 m under four post-processing and registration protocols. Means over three splits.}
\label{tab:protocol}
\footnotesize
\setlength{\tabcolsep}{3.6pt}
\begin{tabular}{@{}lrrrrrrrr@{}}
\toprule
 & \multicolumn{2}{c}{Mesh as extracted} & \multicolumn{2}{c}{+ tie-point alignment} & \multicolumn{2}{c}{+ Taubin smoothing} & \multicolumn{2}{c}{ICP registration} \\
\cmidrule(lr){2-3}\cmidrule(lr){4-5}\cmidrule(lr){6-7}\cmidrule(lr){8-9}
Method & Mean (m) & F@1 m & Mean (m) & F@1 m & Mean (m) & F@1 m & Mean (m) & F@1 m \\
\midrule
\multicolumn{9}{@{}l}{\textit{Itokawa (AMICA, phase 10\textdegree{})}} \\
2DGS & 6.58 & 0.532 & 6.55 & 0.540 & 6.55 & 0.546 & 6.56 & 0.539 \\
GOF & 2.41 & 0.365 & 2.40 & 0.365 & 2.40 & 0.362 & 2.39 & 0.373 \\
PGSR & \underline{1.18} & \underline{0.695} & 0.99 & \underline{0.811} & 0.99 & \underline{0.812} & 0.98 & \underline{0.821} \\
Asteroid-GS & 1.52 & 0.432 & 1.01 & 0.648 & 1.01 & 0.649 & 0.99 & 0.661 \\
Asteroid-GS$^{\dagger}$ & 1.53 & 0.415 & \underline{0.94} & 0.659 & \underline{0.94} & 0.658 & \underline{0.91} & 0.684 \\
\textbf{PCGS (this work)} & \textbf{0.59} & \textbf{0.853} & \textbf{0.59} & \textbf{0.854} & \textbf{0.59} & \textbf{0.853} & \textbf{0.59} & \textbf{0.857} \\
\addlinespace[3pt]
\multicolumn{9}{@{}l}{\textit{Ryugu (ONC-T, phase 19\textdegree{})}} \\
2DGS & 2.51 & \underline{0.388} & 2.51 & 0.410 & 2.50 & 0.416 & 2.48 & 0.421 \\
GOF & 4.20 & 0.156 & 4.21 & 0.177 & 4.20 & 0.162 & 3.38 & 0.255 \\
PGSR & \underline{1.51} & 0.383 & \underline{0.73} & \underline{0.788} & \underline{0.73} & \underline{0.787} & \underline{0.72} & \underline{0.794} \\
Asteroid-GS & 2.41 & 0.239 & 1.27 & 0.524 & 1.27 & 0.525 & 1.26 & 0.530 \\
Asteroid-GS$^{\dagger}$ & 2.43 & 0.242 & 1.29 & 0.529 & 1.29 & 0.530 & 1.27 & 0.539 \\
\textbf{PCGS (this work)} & \textbf{0.62} & \textbf{0.797} & \textbf{0.62} & \textbf{0.798} & \textbf{0.62} & \textbf{0.798} & \textbf{0.61} & \textbf{0.803} \\
\bottomrule
\end{tabular}
\tablefoot{Mesh as extracted: each method's own extraction, registered with the shared scene transform (for PCGS this already contains its offset to the SfM tie points). Tie-point alignment: translation and normal offset fitted to the tie points, no reference used. ICP registration: the mesh after tie-point alignment with a similarity transform (scale, rotation, translation) fitted to the reference by ICP instead of the shared scene transform, which removes any global misregistration of that mesh. $^{\dagger}$Asteroid-GS with the scene extent that its released code derives from the camera positions; the other rows use the object extent. Bold: best; underlined: second best.}
\end{table*}


\begin{table*}[t]
\centering
\caption{Shape and image accuracy on Itokawa and Ryugu. Mean $\pm$ standard deviation over three splits (50 training and 10 test images each).}
\label{tab:main}
\footnotesize
\setlength{\tabcolsep}{2.4pt}
\begin{tabular}{@{}lr@{\,$\pm$\,}lr@{\,$\pm$\,}lr@{\,$\pm$\,}lr@{\,$\pm$\,}lr@{\,$\pm$\,}lr@{\,$\pm$\,}l@{}}
\toprule
 & \multicolumn{10}{c}{Shape against the reference model} & \multicolumn{2}{c}{Images} \\
\cmidrule(lr){2-11}\cmidrule(lr){12-13}
Method & \multicolumn{2}{c}{Mean (m)} & \multicolumn{2}{c}{RMSE (m)} & \multicolumn{2}{c}{Chamfer (m)} & \multicolumn{2}{c}{F@1 m} & \multicolumn{2}{c}{Normal (\textdegree{})} & \multicolumn{2}{c}{PSNR (dB)} \\
\midrule
\multicolumn{13}{@{}l}{\textit{Itokawa (AMICA, phase 10\textdegree{})}} \\
2DGS & 6.55 & 3.41 & 17.96 & 11.62 & 4.63 & 1.92 & 0.546 & 0.035 & 24.1 & 2.4 & 33.6 & 1.0 \\
GOF & 2.40 & 0.12 & 4.66 & 0.45 & 3.01 & 0.24 & 0.362 & 0.004 & 41.5 & 0.4 & \underline{38.0} & 0.6 \\
PGSR & 0.99 & 0.07 & 2.83 & 0.46 & \underline{0.76} & 0.03 & \underline{0.812} & 0.001 & 14.3 & 0.1 & \textbf{39.3} & 0.8 \\
Asteroid-GS & 1.01 & 0.10 & 2.10 & 0.41 & 1.76 & 0.55 & 0.649 & 0.031 & \underline{13.9} & 1.1 & 29.6 & 0.9 \\
Asteroid-GS$^{\dagger}$ & \underline{0.94} & 0.03 & \underline{1.82} & 0.17 & 1.54 & 0.16 & 0.658 & 0.008 & \textbf{12.7} & 0.2 & 29.5 & 0.5 \\
\textbf{PCGS (this work)} & \textbf{0.59} & 0.02 & \textbf{1.02} & 0.15 & \textbf{0.59} & 0.01 & \textbf{0.853} & 0.005 & 14.7 & 0.2 & 36.5 & 0.5 \\
\addlinespace[3pt]
\multicolumn{13}{@{}l}{\textit{Ryugu (ONC-T, phase 19\textdegree{})}} \\
2DGS & 2.50 & 0.08 & 3.80 & 0.17 & 1.96 & 0.07 & 0.416 & 0.010 & 28.0 & 0.4 & \underline{34.9} & 1.8 \\
GOF & 4.20 & 0.08 & 5.98 & 0.15 & 4.50 & 0.11 & 0.162 & 0.003 & 42.4 & 0.3 & 33.5 & 0.0 \\
PGSR & \underline{0.73} & 0.01 & \underline{1.18} & 0.04 & \underline{0.76} & 0.01 & \underline{0.787} & 0.005 & 14.5 & 0.1 & \textbf{35.2} & 0.2 \\
Asteroid-GS & 1.27 & 0.02 & 1.88 & 0.03 & 1.42 & 0.03 & 0.525 & 0.005 & \underline{11.4} & 0.0 & 22.7 & 0.2 \\
Asteroid-GS$^{\dagger}$ & 1.29 & 0.01 & 1.98 & 0.02 & 1.52 & 0.02 & 0.530 & 0.004 & \textbf{11.3} & 0.0 & 22.7 & 0.0 \\
\textbf{PCGS (this work)} & \textbf{0.62} & 0.00 & \textbf{0.89} & 0.01 & \textbf{0.73} & 0.00 & \textbf{0.798} & 0.000 & 12.1 & 0.0 & 31.3 & 0.1 \\
\bottomrule
\end{tabular}
\tablefoot{Distances in metres. Mean, RMSE: unsigned distance from the vertices of the reconstruction to the reference surface; Chamfer: mean of the two directed mean distances; F@1 m: F-score at 1 m; Normal: mean angle between corresponding normals; PSNR: test images after one gain and offset per image. Lower is better except for the F-score and PSNR. Bold: best mean; underlined: second best. $^{\dagger}$Asteroid-GS with the scene extent that its released code derives from the camera positions; the other rows use the object extent.}
\end{table*}


\begin{figure*}
\centering
\includegraphics[width=\textwidth]{\figfile{fig3_numbers}}
\caption{Shape and image accuracy of all methods. (a) Itokawa (Hayabusa AMICA, phase angle 10$^\circ$) and (b) Ryugu (Hayabusa2 ONC-T, 19$^\circ$). Bars give the mean over the three splits, error bars the standard deviation, dots the individual splits, and numbers the mean (bold: best in the panel); bars that exceed the axis are cut and labelled. $^\dagger$Asteroid-GS with the scene extent of its public code. (c) Bennu (OSIRIS-REx MapCam, 26--60$^\circ$), against the OLA shape model, three splits; $^\ddagger$\pcgs\ without the night-side band, first split.}
\label{fig:numbers}
\end{figure*}

\begin{figure*}
\centering
\includegraphics[width=\textwidth]{\figfile{figS1_curves}}
\caption{Accuracy, completeness, and location of the error on Itokawa and Ryugu. (a--d) Share of the reconstructed
surface within a distance $d$ of the reference (accuracy) and of the reference within $d$ of the reconstruction
(completeness), mean over the three splits; the values at $d=1$~m are given for \pcgs\ and the strongest baseline.
The dashed curve is \pcgs\ without priors and without the night-side band. (e, f) Mean completeness
distance per 5$^\circ$ latitude band of the reference; grey shading marks latitudes beyond $\pm60^\circ$.}
\label{fig:curves}
\end{figure*}

%--------------------------------------------------------------------
\section{Results}
\label{sec:results}

\begin{figure*}
\centering
\includegraphics[width=\textwidth,height=0.9\textheight,keepaspectratio]{\figfile{fig2_reconstruction}}
\caption{Reconstructions of Itokawa, Ryugu (19$^\circ$), and Bennu (first split). (a, e, i) A held-out image, labelled with its phase angle, and the reference model, \pcgs, PGSR, and Asteroid-GS rendered from its camera with the same material under low-angle light (Itokawa rotated by 90$^\circ$); yellow boxes mark the windows. (b, f, j) The window, placed automatically where the reference model has the most local relief, without reference to any reconstruction. (c, g, k) Distribution of the signed distance $d$ from the reconstructed surface to the reference (positive outside). (d, h, l) $d$ on the surface seen from the same camera, clipped at $\pm3$~m, with the mean $|d|$ and the RMSE of the vertices. Asteroid-GS uses the object extent on Itokawa and Ryugu and the scene extent of its public code on Bennu.}
\label{fig:recon}
\end{figure*}

\subsection{Accuracy against the image-based references}
\label{sec:accuracy}

On Itokawa and on the 19$^\circ$ Ryugu image set, \pcgs\ is the most accurate method by every distance measure
(Table~\ref{tab:main} and Fig.~\ref{fig:numbers}a,b). It has the lowest mean distance, RMSE, and Chamfer distance and
the highest F-score at 1~m on both bodies. Its mean distance is 0.59~m on Itokawa and 0.62~m on Ryugu, which is 37\%
lower than that of the best competing method on Itokawa, Asteroid-GS$^\dagger$, and 15\% lower than that of the best
competing method on Ryugu, PGSR. The RMSE is lower by 44\% and 24\%, respectively. The individual splits scatter little
around these means, so the ranking does not depend on the image selection. 2DGS and GOF are considerably less accurate,
with mean distances of 2.4--6.6~m on both bodies.

The lower errors are not obtained at the expense of completeness. Both the accuracy and the completeness at 1~m exceed
those of every other method (Table~\ref{tab:shape_full}), and the cumulative distributions show the largest lead at
sub-metre distances (Fig.~\ref{fig:curves}a--d). On Ryugu, \pcgs\ and PGSR perform similarly over most of the surface:
their Chamfer distances differ by only 0.03~m, and PGSR has a slightly higher F-score at 2~m. Here, the advantage of
\pcgs\ lies mainly in the absence of large errors, which is why it is larger in RMSE than in the other measures.

\begin{table*}[t]
\centering
\caption{Shape metrics in full. Mean $\pm$ s.d. over three splits; protocol of Table~\ref{tab:main}.}
\label{tab:shape_full}
\scriptsize
\setlength{\tabcolsep}{1.7pt}
\begin{tabular}{@{}lr@{\,$\pm$\,}lr@{\,$\pm$\,}lr@{\,$\pm$\,}lr@{\,$\pm$\,}lr@{\,$\pm$\,}lr@{\,$\pm$\,}lr@{\,$\pm$\,}lr@{\,$\pm$\,}lr@{\,$\pm$\,}lr@{\,$\pm$\,}l@{}}
\toprule
Method & \multicolumn{2}{c}{Mean (m) $\downarrow$} & \multicolumn{2}{c}{RMSE (m) $\downarrow$} & \multicolumn{2}{c}{Std (m) $\downarrow$} & \multicolumn{2}{c}{Chamfer (m) $\downarrow$} & \multicolumn{2}{c}{Acc. (\%) $\uparrow$} & \multicolumn{2}{c}{Compl. (\%) $\uparrow$} & \multicolumn{2}{c}{F@0.5 m $\uparrow$} & \multicolumn{2}{c}{F@1 m $\uparrow$} & \multicolumn{2}{c}{F@2 m $\uparrow$} & \multicolumn{2}{c}{Normal (\textdegree{}) $\downarrow$} \\
\midrule
\multicolumn{21}{@{}l}{\textit{Itokawa (AMICA, phase 10\textdegree{})}} \\
2DGS & 6.55 & 3.41 & 17.96 & 11.62 & 16.70 & 11.16 & 4.63 & 1.92 & 50.3 & 4.7 & 59.8 & 1.8 & 0.338 & 0.024 & 0.546 & 0.035 & 0.711 & 0.037 & 24.1 & 2.4 \\
GOF & 2.40 & 0.12 & 4.66 & 0.45 & 4.00 & 0.46 & 3.01 & 0.24 & 28.2 & 0.2 & 50.2 & 1.0 & 0.196 & 0.004 & 0.362 & 0.004 & 0.593 & 0.006 & 41.5 & 0.4 \\
PGSR & 0.99 & 0.07 & 2.83 & 0.46 & 2.66 & 0.46 & \underline{0.76} & 0.03 & \underline{80.7} & 0.2 & \underline{81.8} & 0.2 & \underline{0.526} & 0.001 & \underline{0.812} & 0.001 & \underline{0.958} & 0.001 & 14.3 & 0.1 \\
Asteroid-GS & 1.01 & 0.10 & 2.10 & 0.41 & 1.83 & 0.41 & 1.76 & 0.55 & 63.5 & 4.0 & 66.5 & 2.1 & 0.376 & 0.022 & 0.649 & 0.031 & 0.835 & 0.028 & \underline{13.9} & 1.1 \\
Asteroid-GS$^{\dagger}$ & \underline{0.94} & 0.03 & \underline{1.82} & 0.17 & \underline{1.55} & 0.19 & 1.54 & 0.16 & 65.9 & 0.8 & 65.8 & 0.8 & 0.378 & 0.003 & 0.658 & 0.008 & 0.848 & 0.007 & \textbf{12.7} & 0.2 \\
Mip-Splatting$^{\ddagger}$ & \multicolumn{2}{c}{2.48} & \multicolumn{2}{c}{3.20} & \multicolumn{2}{c}{2.01} & \multicolumn{2}{c}{73.26} & \multicolumn{2}{c}{31.5} & \multicolumn{2}{c}{10.5} & \multicolumn{2}{c}{0.086} & \multicolumn{2}{c}{0.158} & \multicolumn{2}{c}{0.247} & \multicolumn{2}{c}{50.0} \\
\textbf{PCGS (this work)} & \textbf{0.59} & 0.02 & \textbf{1.02} & 0.15 & \textbf{0.83} & 0.17 & \textbf{0.59} & 0.01 & \textbf{84.7} & 0.7 & \textbf{85.9} & 0.4 & \textbf{0.579} & 0.005 & \textbf{0.853} & 0.005 & \textbf{0.969} & 0.002 & 14.7 & 0.2 \\
\addlinespace[3pt]
\multicolumn{21}{@{}l}{\textit{Ryugu (ONC-T, phase 19\textdegree{})}} \\
2DGS & 2.50 & 0.08 & 3.80 & 0.17 & 2.86 & 0.16 & 1.96 & 0.07 & 37.3 & 1.1 & 47.1 & 0.8 & 0.222 & 0.006 & 0.416 & 0.010 & 0.679 & 0.008 & 28.0 & 0.4 \\
GOF & 4.20 & 0.08 & 5.98 & 0.15 & 4.25 & 0.14 & 4.50 & 0.11 & 12.9 & 0.3 & 22.0 & 0.3 & 0.082 & 0.002 & 0.162 & 0.003 & 0.319 & 0.006 & 42.4 & 0.3 \\
PGSR & \underline{0.73} & 0.01 & \underline{1.18} & 0.04 & \underline{0.93} & 0.04 & \underline{0.76} & 0.01 & \underline{79.1} & 0.5 & \underline{78.2} & 0.5 & \underline{0.499} & 0.007 & \underline{0.787} & 0.005 & \textbf{0.949} & 0.001 & 14.5 & 0.1 \\
Asteroid-GS & 1.27 & 0.02 & 1.88 & 0.03 & 1.39 & 0.03 & 1.42 & 0.03 & 53.2 & 0.6 & 51.8 & 0.5 & 0.284 & 0.004 & 0.525 & 0.005 & 0.797 & 0.005 & \underline{11.4} & 0.0 \\
Asteroid-GS$^{\dagger}$ & 1.29 & 0.01 & 1.98 & 0.02 & 1.51 & 0.02 & 1.52 & 0.02 & 53.7 & 0.4 & 52.3 & 0.4 & 0.284 & 0.002 & 0.530 & 0.004 & 0.794 & 0.004 & \textbf{11.3} & 0.0 \\
\textbf{PCGS (this work)} & \textbf{0.62} & 0.00 & \textbf{0.89} & 0.01 & \textbf{0.65} & 0.01 & \textbf{0.73} & 0.00 & \textbf{80.2} & 0.0 & \textbf{79.4} & 0.0 & \textbf{0.516} & 0.001 & \textbf{0.798} & 0.000 & \underline{0.948} & 0.001 & 12.1 & 0.0 \\
\bottomrule
\end{tabular}
\tablefoot{Acc. and Compl.: shares of reconstructed and reference samples within 1 m of the other surface; F@d their harmonic mean at distance d. Chamfer: mean of the two directed mean distances. Std: standard deviation of the vertex distances. $^{\dagger}$Asteroid-GS with the scene extent of its public repository (camera ring); the other rows use the object extent, as all methods do. $^{\ddagger}$First split only; mesh by TSDF fusion of its rendered depth. Bold: best; underlined: second best.}
\end{table*}


Because the meshes of \pcgs\ are closed, their volumes can be compared directly with those of the references. At the
registered scale, they agree to within 0.3\% on Itokawa and 0.01\% on Ryugu (Table~\ref{tab:volume}). The normal error
is the only measure in which \pcgs\ does not lead. Asteroid-GS has the lowest normal error on both bodies, and on
Itokawa PGSR is also slightly better than \pcgs. Smooth surfaces carry little high-frequency normal error, which
favours Asteroid-GS. Part of the difference may also reflect noise in the surfel normals at the pixel scale, where the
shading constrains the orientation only up to the image resolution.

\begin{table}
\caption{Volume error of the \pcgs\ meshes relative to the reference models.}
\label{tab:volume}
\centering
\footnotesize\setlength{\tabcolsep}{4pt}
\begin{tabular}{@{}lccc@{}}
\toprule
 & & \multicolumn{2}{c}{Volume error (\%)} \\
\cmidrule(lr){3-4}
Image set & Scale & registered & SPICE scale \\
\midrule
Ryugu, 7$^\circ$  & 0.9993 & $+0.11$ & $+0.32$ \\
Itokawa, 10$^\circ$ & 1.0105 & $-0.19$, $-0.27$, $-0.25$ & $-3.27$, $-3.35$, $-3.33$ \\
Ryugu, 19$^\circ$ & 0.9985 & $+0.01$, $+0.01$, $+0.00$ & $+0.46$, $+0.46$, $+0.45$ \\
Ryugu, 40$^\circ$ & 0.9969 & $+0.06$, $+0.05$ & $+1.00$, $+0.99$ \\
Bennu, 26--60$^\circ$ & 0.9981 & $+0.25$, $+0.29$, $+0.32$ & $+0.82$, $+0.86$, $+0.89$ \\
\bottomrule
\end{tabular}
\tablefoot{One value per available split. Scale: that of the similarity transform from the camera frame, which has
the SPICE scale, to the reference frame (Sect.~\ref{sec:reference}). Registered: the volume after this transform;
SPICE scale: the volume in the camera frame. Reference volumes: 0.017734~km$^3$ (Itokawa), 0.37743~km$^3$ (Ryugu),
and 0.061316~km$^3$ (Bennu). Some meshes contain up to five non-manifold edges among about $10^7$ and no boundary
edges; their volumes are computed from the closed surface. Of the baselines, only GOF produces closed meshes
consistently, with registered volume errors of $-1.4$ to $-3.0$\% on Itokawa and $+0.5$ to $+0.6$\% on Ryugu.}
\end{table}

The reconstructions in Fig.~\ref{fig:recon}a--h show where these differences arise. All methods recover the global
shapes, that is, the elongated, two-lobed shape of Itokawa and the spinning-top shape of Ryugu with its equatorial
ridge; they differ mainly at the scale of individual boulders. In the windows of greatest relief
(Fig.~\ref{fig:recon}b,f), \pcgs\ reproduces the boulder-strewn terrain of the references, PGSR shows similar relief with
streak artefacts, and Asteroid-GS smooths the surface so that only the largest blocks remain. Accordingly, the signed
distance of \pcgs\ is narrowly distributed around zero, whereas that of Asteroid-GS is broad and has a long tail of
surface outside the reference (Fig.~\ref{fig:recon}c,g). The largest errors of \pcgs\ occur in a few concave regions of
Itokawa that are poorly covered by the images (Fig.~\ref{fig:recon}d). PGSR has errors in the same regions, and
Asteroid-GS has extended areas lying 1--3~m outside the reference on both bodies (Fig.~\ref{fig:recon}d,h). Overall,
\pcgs\ recovers more of the boulder-scale relief than Asteroid-GS and avoids the large errors that dominate the RMSE of
PGSR.

\subsection{Validation against an image-independent reference}
\label{sec:bennu}

Because the Itokawa and Ryugu references were derived from images, they could favour methods that share their
systematic errors. Bennu provides an independent test: its reference was derived from laser ranges alone, and its images
span phase angles of 26--60$^\circ$, so much of the body is unlit in every image. The differences between the methods
are largest on this body (Table~\ref{tab:bennu} and Fig.~\ref{fig:numbers}c). \pcgs\ reaches a mean distance of
$0.62\pm0.06$~m, as small as against the image-based references, and an F-score at 1~m of $0.812\pm0.011$. PGSR and
Asteroid-GS$^\dagger$ both reach mean distances of 3.78~m, with large variations between splits, and 2DGS fails on this
image set. The RMSE and the normal error of \pcgs\ are higher than on the other bodies. This is partly due to a region of
large error described below and partly because the reference resolves the boulder-covered surface at 0.88~m, whereas
the images have pixel scales of up to 2.6~m.

\begin{table*}[t]
\centering
\caption{Shape and image accuracy on Bennu against the laser-altimeter model. Mean $\pm$ standard deviation over three splits.}
\label{tab:bennu}
\footnotesize
\setlength{\tabcolsep}{2.4pt}
\begin{tabular}{@{}lr@{\,$\pm$\,}lr@{\,$\pm$\,}lr@{\,$\pm$\,}lr@{\,$\pm$\,}lr@{\,$\pm$\,}lr@{\,$\pm$\,}l@{}}
\toprule
 & \multicolumn{10}{c}{Shape against the OLA model} & \multicolumn{2}{c}{Images} \\
\cmidrule(lr){2-11}\cmidrule(lr){12-13}
Method & \multicolumn{2}{c}{Mean (m) $\downarrow$} & \multicolumn{2}{c}{RMSE (m) $\downarrow$} & \multicolumn{2}{c}{Chamfer (m) $\downarrow$} & \multicolumn{2}{c}{F@1 m $\uparrow$} & \multicolumn{2}{c}{Normal (\textdegree{}) $\downarrow$} & \multicolumn{2}{c}{PSNR (dB) $\uparrow$} \\
\midrule
2DGS & 26.77 & 1.59 & 41.64 & 2.69 & 26.85 & 0.33 & 0.138 & 0.023 & 53.3 & 0.8 & 20.1 & 0.6 \\
PGSR & \underline{3.78} & 0.85 & 11.26 & 3.39 & \underline{3.84} & 0.81 & \underline{0.689} & 0.049 & 38.0 & 1.3 & \underline{20.3} & 0.3 \\
Asteroid-GS$^{\dagger}$ & \underline{3.78} & 1.75 & \underline{8.69} & 5.00 & 5.94 & 3.47 & 0.368 & 0.035 & \underline{33.1} & 1.3 & 17.6 & 1.0 \\
\textbf{PCGS (this work)} & \textbf{0.62} & 0.06 & \textbf{1.34} & 0.14 & \textbf{0.75} & 0.04 & \textbf{0.812} & 0.011 & \textbf{32.8} & 0.3 & \textbf{22.2} & 0.5 \\
\bottomrule
\end{tabular}
\tablefoot{Distances in metres. Mean, RMSE: vertices of the reconstruction to the reference surface; Chamfer: mean of the two directed mean distances; F@1 m: F-score at 1 m; Normal: mean angle between corresponding normals; PSNR: held-out images after one gain and offset per image. Bold: best; underlined: second best. $^{\dagger}$Asteroid-GS with the scene extent of its public code (the other rows use the object extent). GOF exceeded the available GPU memory on Bennu.}
\end{table*}


The reconstructions in Fig.~\ref{fig:recon}i--l show how the baselines fail. PGSR leaves holes and displaced patches,
and Asteroid-GS produces a smooth body that lies far inside the reference over a large region: 17.3\% of the PGSR
vertices lie more than 3~m outside the reference, and 11.3\% of the Asteroid-GS vertices more than 3~m inside it. The
mesh of \pcgs\ is complete, reproduces the equatorial ridge and the large boulders, and has a signed distance centred on
zero. Its largest errors form a single contiguous region, containing 1.4\% of the vertices, in which the surface lies
more than 3~m outside the reference. This region may consist of surface that grew into the night-side band, where
neither the silhouette loss nor the unlit pixels penalise it. At the scale of the window (Fig.~\ref{fig:recon}j), \pcgs\
recovers a few large blocks but smooths the dense field of small boulders resolved by the laser altimeter. The window is
resolved at 0.67~m per pixel in this image, but only at up to 2.6~m per pixel in other training images.

Most of the difference from the baselines results from the treatment of the night side. Without the night-side band,
the mean distance of \pcgs\ on the first split is 4.18~m, larger than that of PGSR and Asteroid-GS$^\dagger$ on the same
split, 4.06 and 2.85~m (Table~\ref{tab:phase}). Figure~\ref{fig:night} illustrates the mechanism. In a test image at a
phase angle of 45$^\circ$, the outline of the reference model extends well beyond the foreground mask into the band
defined by Eq.~(\ref{eq:night}) (Fig.~\ref{fig:night}a). Without the band, the silhouette loss treats this region as
empty space: the surface is removed where the body is dark in most training views and bulges outwards elsewhere, and
16\% of the vertices lie more than 6~m from the reference (Fig.~\ref{fig:night}b,d). With the band, this fraction falls
to 1\%, and the signed distance is centred on zero. At 40$^\circ$ on Ryugu, the band is narrower but still covers a
crescent several tens of pixels wide (Fig.~\ref{fig:night}c). The baselines show removed and displaced regions of the
same kind (Fig.~\ref{fig:recon}l), but they were trained without the band. The comparison on Bennu thus mainly shows
that the night side must be treated explicitly at high phase angle; it does not isolate the contribution of the
photometric model.

\begin{figure*}
\centering
\includegraphics[width=\textwidth]{\figfile{fig_night}}
\caption{The night side at high phase angle (first split). (a) A held-out Bennu image (OSIRIS-REx MapCam, phase angle
45$^\circ$) with the foreground mask (yellow), the night-side band of Eq.~(\ref{eq:night}) (blue), and the outline of the
OLA reference projected into the same view (white, dashed), followed by the reference and \pcgs\ without and with the
band, rendered from the same camera. (b) Distribution of the signed distance from the mesh vertices to the reference,
with the share of vertices farther than 6~m. (c) The band for a held-out Ryugu image at 40$^\circ$. (d) Signed
distance on the surface without and with the band, clipped at $\pm3$~m.}
\label{fig:night}
\end{figure*}

\subsection{Spatial distribution of the error}
\label{sec:maps}

The errors of every method are concentrated where the observation geometry constrains the surface poorly. On Itokawa,
the completeness distance of \pcgs\ varies little with latitude, whereas PGSR has large errors at southern latitudes and
Asteroid-GS at mid-latitudes (Fig.~\ref{fig:curves}e). On Ryugu, all methods are least accurate near the poles, which the
near-equatorial images observe only at grazing angles (Fig.~\ref{fig:curves}f). The completeness distance of \pcgs\ and
PGSR increases from 0.5--0.7~m at low latitudes to 2.5--3.7~m near the poles, and that of Asteroid-GS is higher at
nearly all latitudes.

The maps of the completeness distance over the whole surface confirm this picture (Fig.~\ref{fig:maps}). On Itokawa and
Ryugu, \pcgs\ and PGSR have small errors over most of the surface, apart from a few isolated patches and, on Ryugu, the
two poles, where all methods have large errors. Asteroid-GS has large contiguous regions with errors above 2.5~m on all
three bodies (Fig.~\ref{fig:maps}c,f,i). On Bennu, PGSR and Asteroid-GS miss large parts of the body over a range of
longitudes, with mean completeness distances of 5.7 and 5.2~m, whereas the errors of \pcgs, with a mean of 0.69~m, are
concentrated in a narrow band along one meridian and in a few patches (Fig.~\ref{fig:maps}g). The large errors of
\pcgs\ are thus localised, whereas the baselines leave extended regions poorly reconstructed.

\begin{figure*}
\centering
\includegraphics[width=\textwidth]{\figfile{fig_maps}}
\caption{Completeness distance (from the reference surface to the reconstruction) over the whole surface of Itokawa (a--c), Ryugu at 19$^\circ$ (d--f), and Bennu (g--i), first split, in Mollweide projections of the body-fixed frame (longitude 0 at the centre, grid lines every 60$^\circ$ in longitude and 30$^\circ$ in latitude). Rows: \pcgs, PGSR, and Asteroid-GS (on Bennu with the scene extent of its public code). Values are averaged in $2^\circ$ bins and smoothed with a Gaussian kernel of 1.2 bins, empty bins excluded; the mean is given in each panel.}
\label{fig:maps}
\end{figure*}

\subsection{Dependence on the phase angle}
\label{sec:phase}

The five image sets span phase angles from 7$^\circ$ to 60$^\circ$, and three of them show the same body
(Table~\ref{tab:phase}). The mean distance of \pcgs\ remains between 0.59 and 0.75~m over the whole range and is lower
than that of the baselines on every image set on which they were run. On Ryugu, it is lowest at 19$^\circ$ and higher
at 7$^\circ$ and 40$^\circ$. Weaker shading at 7$^\circ$ and a larger night side at 40$^\circ$ may both contribute, as
may the use of the 19$^\circ$ image set during development (Sect.~\ref{sec:images}).

\begin{table*}[t]
\centering
\caption{Shape accuracy on the five image sets, ordered by phase angle. First split; mean distance (m) and F-score at 1 m.}
\label{tab:phase}
\footnotesize
\setlength{\tabcolsep}{2.0pt}
\begin{tabular}{@{}lrrrrrrrrrr@{}}
\toprule
 & \multicolumn{2}{c}{Ryugu, 7\textdegree{}} & \multicolumn{2}{c}{Itokawa, 10\textdegree{}} & \multicolumn{2}{c}{Ryugu, 19\textdegree{}} & \multicolumn{2}{c}{Ryugu, 40\textdegree{}} & \multicolumn{2}{c}{Bennu, 26--60\textdegree{}} \\
\cmidrule(lr){2-3}\cmidrule(lr){4-5}\cmidrule(lr){6-7}\cmidrule(lr){8-9}\cmidrule(lr){10-11}
Method & Mean (m) & F@1 m & Mean (m) & F@1 m & Mean (m) & F@1 m & Mean (m) & F@1 m & Mean (m) & F@1 m \\
\midrule
PGSR & 0.84 & \textbf{0.730} & 0.90 & 0.812 & 0.72 & 0.789 & 0.92 & 0.711 & 4.06 & \underline{0.642} \\
Asteroid-GS & 1.45 & 0.491 & 1.10 & 0.621 & 1.29 & 0.519 & 1.25 & 0.527 & 2.85 & 0.408 \\
PCGS without phase switch & \textbf{0.74} & \underline{0.720} & 0.68 & 0.841 & \underline{0.64} & 0.779 & \underline{0.68} & \underline{0.743} & \underline{2.44} & 0.606 \\
PCGS without night-side band & 0.76 & 0.708 & \textbf{0.58} & \underline{0.846} & \textbf{0.62} & \underline{0.791} & 0.71 & 0.709 & 4.18 & 0.444 \\
\textbf{PCGS (this work)} & \underline{0.75} & 0.712 & \underline{0.60} & \textbf{0.849} & \textbf{0.62} & \textbf{0.798} & \textbf{0.66} & \textbf{0.757} & \textbf{0.59} & \textbf{0.811} \\
\bottomrule
\end{tabular}
\tablefoot{Phase switch: albedo prior below 15\textdegree{}, cast shadows at and above 15\textdegree{}. The Bennu entry of Asteroid-GS is Asteroid-GS$^{\dagger}$ (scene extent of its released code). --: not run. Lower mean distance and higher F-score are better. Bold: best; underlined: second best.}
\end{table*}


The three Ryugu image sets, acquired months apart under different illumination, yield consistent shapes
(Fig.~\ref{fig:ryuguphase}). Seen from the same camera, the three meshes are nearly indistinguishable at the scale of the
whole body and of the large boulders (Fig.~\ref{fig:ryuguphase}a,b). Their differences are systematic but small: the
signed distance is shifted slightly outwards at 7$^\circ$ and slightly inwards at 40$^\circ$, and it is centred on the
reference at 19$^\circ$ (Fig.~\ref{fig:ryuguphase}c). Asteroid-GS, trained on the 40$^\circ$ images, has a broad error
distribution with deviations of several metres in both directions.

The learned partition parameter $L$ of the lunar-Lambert function also varies with the phase angle. On Ryugu, it
decreases from 0.78 at 7$^\circ$ to 0.62 at 19$^\circ$ and 0.37 at 40$^\circ$, with differences between splits below
0.01, and its value for Itokawa at 10$^\circ$, 0.64, lies between the first two. \citet{McEwen1991} found the same trend
for the Moon. On Bennu, $L$ is about 0.1; because Eq.~(\ref{eq:ll}) contains no phase function and the Bennu images span
26--60$^\circ$, this value probably also absorbs the phase dependence of the brightness and should not be interpreted as
a photometric measurement.

\begin{figure*}
\centering
\includegraphics[width=\textwidth]{\figfile{fig_ryugu_phase}}
\caption{Ryugu reconstructed from the three image sets at phase angles of 7$^\circ$, 19$^\circ$, and 40$^\circ$
(first split), and by Asteroid-GS from the 40$^\circ$ image set, all registered to the reference frame and rendered from
the camera of the 19$^\circ$ held-out image of Fig.~\ref{fig:recon}. (a) Full view. (b) The window of
Fig.~\ref{fig:recon}f. (c) Distribution of the signed distance to the reference and the signed distance on the
surface, clipped at $\pm3$~m, with the mean $|d|$ and the RMSE.}
\label{fig:ryuguphase}
\end{figure*}

\subsection{Ablation study}
\label{sec:ablation}

We now examine the contribution of each component of \pcgs: the phase-dependent prior, the night-side band, the
photometric rendering, and the surface extraction.

\subsubsection{Phase-dependent prior}
\label{sec:rule}

The phase-dependent prior addresses the two conditions in which photoclinometry from a single image set is weakest: weak
shading at low phase angle and frequent cast shadows at high phase angle. To isolate its effect,
Fig.~\ref{fig:priors} and Table~\ref{tab:priors} compare models with and without the prior selected for each phase
angle, both trained without the night-side band. The albedo prior has its largest effect on Itokawa, where it reduces
the RMSE of the first split from 1.46 to 0.92~m. Without the prior, part of the weak shading at 10$^\circ$ is attributed
to albedo variations instead of slopes, and larger areas of the surface are displaced (Fig.~\ref{fig:priors}c). On
Ryugu at 19$^\circ$, cast shadows improve the RMSE, the F-score, and the normal error on each of the three splits
(Fig.~\ref{fig:priors}b) and reduce the mean distance from 0.65 to 0.62~m (Fig.~\ref{fig:priors}d). On Bennu, they
reduce the RMSE by 63\%, although both models remain far from the reference without the band.

The prior brings no improvement at 7$^\circ$ and 40$^\circ$. At 7$^\circ$, the albedo prior slightly lowers the F-score,
from 0.719 to 0.708; at 40$^\circ$, cast shadows reduce the normal error by 0.7$^\circ$ but do not change the distances.
These differences amount to about 0.01~m in mean distance and are based on one or two splits. Using both priors
together is worse than using cast shadows alone at 19$^\circ$ and 40$^\circ$, possibly because the albedo prior then
suppresses the brightness contrast that the shadows explain. As no image set lies between 10.4$^\circ$ and
18.5$^\circ$, the data constrain the threshold only to this interval.

\begin{table*}[t]
\centering
\caption{Effect of the albedo prior and cast shadows. First split; models without the night-side band, and the final model.}
\label{tab:priors}
\footnotesize
\setlength{\tabcolsep}{4.5pt}
\begin{tabular}{@{}lcrrrrr@{}}
\toprule
Method & Switch & Mean (m) & RMSE (m) & F@1 m & Normal (\textdegree{}) & PSNR (dB) \\
\midrule
\multicolumn{7}{@{}l}{\textit{Itokawa, 10\textdegree{}}} \\
Shading only &  & 0.66 & 1.46 & 0.827 & \underline{14.6} & \textbf{38.1} \\
+ albedo prior & \checkmark & \textbf{0.58} & \textbf{0.92} & \underline{0.846} & \textbf{14.4} & 36.8 \\
PCGS (final model) &  & \underline{0.60} & \underline{1.07} & \textbf{0.849} & 14.8 & \underline{37.2} \\
\addlinespace[3pt]
\multicolumn{7}{@{}l}{\textit{Ryugu, 7\textdegree{}}} \\
Shading only &  & \textbf{0.74} & \textbf{1.04} & \textbf{0.719} & \textbf{12.6} & \textbf{34.8} \\
+ albedo prior & \checkmark & 0.76 & \textbf{1.04} & 0.708 & \underline{12.7} & \underline{33.7} \\
PCGS (final model) &  & \underline{0.75} & \underline{1.05} & \underline{0.712} & 13.0 & 33.5 \\
\addlinespace[3pt]
\multicolumn{7}{@{}l}{\textit{Ryugu, 19\textdegree{}}} \\
Shading only &  & 0.65 & 0.95 & 0.770 & 12.4 & \textbf{32.5} \\
+ albedo prior &  & 0.66 & 0.96 & 0.767 & 13.1 & 31.0 \\
+ cast shadows & \checkmark & \textbf{0.62} & \underline{0.90} & \underline{0.791} & \textbf{11.9} & 30.0 \\
+ albedo prior + cast shadows &  & \underline{0.63} & 0.93 & 0.785 & 12.4 & 28.6 \\
PCGS (final model) &  & \textbf{0.62} & \textbf{0.89} & \textbf{0.798} & \underline{12.1} & \underline{31.4} \\
\addlinespace[3pt]
\multicolumn{7}{@{}l}{\textit{Ryugu, 40\textdegree{}}} \\
Shading only &  & \underline{0.71} & \underline{1.06} & \underline{0.719} & 13.2 & \textbf{33.9} \\
+ cast shadows & \checkmark & \underline{0.71} & 1.07 & 0.709 & \textbf{12.5} & 28.1 \\
+ albedo prior + cast shadows &  & 0.77 & 1.24 & 0.676 & 13.0 & 26.6 \\
PCGS (final model) &  & \textbf{0.66} & \textbf{0.95} & \textbf{0.757} & \underline{12.8} & \underline{33.0} \\
\addlinespace[3pt]
\multicolumn{7}{@{}l}{\textit{Bennu, 26--60\textdegree{}}} \\
Shading only &  & 14.75 & 27.49 & 0.271 & 38.6 & \underline{20.3} \\
+ cast shadows & \checkmark & \underline{4.18} & \underline{10.28} & \underline{0.444} & \underline{34.8} & 19.8 \\
PCGS (final model) &  & \textbf{0.59} & \textbf{1.21} & \textbf{0.811} & \textbf{32.9} & \textbf{22.7} \\
\bottomrule
\end{tabular}
\tablefoot{Switch: the variant selected by the phase switch; the final model adds the night-side band to it. Shading only: photometric rendering without albedo prior and cast shadows. Lower is better except for the F-score and PSNR. Bold: best; underlined: second best (image sets with three or more rows).}
\end{table*}


\begin{figure*}
\centering
\includegraphics[width=\textwidth]{\figfile{fig5_priors}}
\caption{The phase rule; all models without the night-side band. (a) One held-out image of each image set except Bennu, ordered by phase angle; the bar below names the prior that the rule selects. (b) RMSE, F-score at 1~m, and normal error without a prior (open circles) and with the prior selected by the rule (filled) for every available split; vermilion lines mark an improvement and grey lines a deterioration; Bennu lies outside the axes. (c, d) Signed distance $d$ without and with the prior on Itokawa (albedo total variation) and on Ryugu at 19$^\circ$ (cast shadows), first split, clipped at $\pm3$~m.}
\label{fig:priors}
\end{figure*}

\subsubsection{Night-side band}

The band has little effect at low phase angle, and its effect increases with the phase angle. On Itokawa, it changes the
mean distance by less than the scatter between splits, from 0.60 to 0.59~m, although it increases the RMSE of the first
split from 0.92 to 1.07~m. On the first split of Ryugu at 40$^\circ$, it reduces the mean distance from 0.71 to 0.66~m,
and on the first split of Bennu from 4.18 to 0.59~m.

\subsubsection{Photometric rendering and multi-view consistency}

Shading and multi-view consistency complement each other, as anticipated in Sect.~\ref{sec:losses}.
Table~\ref{tab:development} lists models in which the components are added one at a time on the first Itokawa split.
These models were evaluated with an earlier protocol, with $3\times10^5$ samples and without alignment or smoothing, so
only the differences between rows are meaningful. The median depth and the silhouette loss reduce the mean distance of
the default 2DGS model from 10.84 to 2.37~m. With the view-dependent colour of 2DGS, multi-view consistency then reduces
it to 1.24~m. Replacing this colour by the shaded albedo reduces it further to 0.82~m and increases the fraction of
vertices within 1~m from 73.7\% to 76.5\%. Shading without multi-view consistency, by contrast, reduces the mean distance
only to 2.15~m and lowers the fraction within 1~m from 59.2\% to 55.7\%. The photometric function itself is also
important: a fixed view-independent albedo without it gives 1.45~m on Itokawa and fails on Ryugu at 40$^\circ$, with
errors of 23--27~m, whereas the shaded model reaches 1.48~m.

\begin{table}
\caption{Development of the model on Itokawa (first split).}
\label{tab:development}
\centering
\footnotesize\setlength{\tabcolsep}{2.5pt}
\begin{tabular}{@{}lccc@{}}
\toprule
Model & Mean / RMSE (m) & $<$1~m & Normal ($^\circ$) \\
\midrule
2DGS, default settings              & 10.84 / 41.59 & 32.5\% & 42.0 \\
+ median depth                      & 2.71 / 6.49   & 56.8\% & 28.0 \\
+ silhouette (A)                    & 2.37 / 5.25   & 59.2\% & 29.8 \\
A + shading                         & 2.15 / 4.83   & 55.7\% & 27.0 \\
A + multi-view NCC                  & 1.24 / 3.55   & 73.7\% & 21.7 \\
A + shading + NCC (B)               & 0.82 / 1.97   & 76.5\% & 18.8 \\
B without shading$^{(a)}$           & 1.45 / 4.83   & 73.9\% & 24.2 \\
B + pixel convention                & 0.71 / 1.94   & 84.0\% & 17.3 \\
\bottomrule
\end{tabular}
\tablefoot{$<$1~m: share of vertices within 1~m of the reference. All rows use the consistent fusion of
Sect.~\ref{sec:extraction}, without priors, hole fairing, or offset; rows without shading keep the view-dependent colour of 2DGS. $^{(a)}$Fixed view-independent albedo without
the photometric function. On Ryugu at 40$^\circ$, model B reached 1.48~m and the model without shading failed
(23--27~m).}
\end{table}

\subsubsection{Surface extraction and pixel convention}

The surface extraction matters at least as much as the training. With the same trained surfels, a signed-distance
fusion of all rendered depths gives mean distances of 7--10~m, compared with 0.8--1.2~m for the fusion of consistent
depth. Correcting the pixel convention reduces the mean distance on Ryugu from 1.35 to 0.76~m, and the offset to the tie
points reduces it further to 0.66~m (Table~\ref{tab:pipeline}). On Itokawa, the fairing mainly reduces the RMSE,
from 1.71 to 1.29~m, because it corrects a few concave regions where the Poisson surface was pulled inwards.

\begin{table}
\caption{Steps of the extraction and of the pixel convention (no priors).}
\label{tab:pipeline}
\centering
\footnotesize\setlength{\tabcolsep}{2pt}
\begin{tabular}{@{}lcc@{}}
\toprule
Step & Itokawa & Ryugu \\
\midrule
Initial extraction                         & 0.78 / 1.56 / 76.9\% & 1.35 / 1.66 / 40.8\% \\
+ pixel conv.\ (extraction)                & 0.68 / 1.52 / 83.3\% & 0.94 / 1.24 / 62.9\% \\
+ pixel conv.\ (training)                  & 0.69 / 1.71 / 84.3\% & 0.76 / 1.08 / 73.5\% \\
+ fairing                             & 0.65 / 1.29 / 84.3\% & 0.76 / 1.06 / 73.4\% \\
+ offset to the SfM tie points             & 0.64 / 1.29 / 84.8\% & 0.66 / 0.95 / 80.3\% \\
\bottomrule
\end{tabular}
\tablefoot{Mean distance (m) / RMSE (m) / share of vertices within 1~m, averaged over three splits; Ryugu at
19$^\circ$.}
\end{table}

Overall, no single component accounts for the accuracy of \pcgs. The photometric rendering is effective only in
combination with multi-view consistency, the night-side band is decisive only at high phase angle, and the
consistency-filtered extraction is needed to convert the trained surfels into an accurate mesh.

\begin{figure*}
\centering
\includegraphics[width=\textwidth]{\figfile{fig4_photometry}}
\caption{Held-out images. (a) A held-out Ryugu image (Hayabusa2 ONC-T, phase angle 19$^\circ$) and the renders of
\pcgs, PGSR, and Asteroid-GS for its camera after the per-image gain and offset, with the PSNR and SSIM of this
image; the yellow box marks the window in (b). (b) The window. (c) PSNR of every held-out image of the three splits
(30 per method and body); black lines mark the median. (d) Render minus image on the body, with its RMS. The render
of \pcgs\ contains no cast shadows; its PSNR is below that of PGSR, whereas its shape error is smaller.}
\label{fig:nvs}
\end{figure*}

\subsection{Image fidelity and shape accuracy}
\label{sec:nvs}

An accurate reproduction of test images does not imply an accurate shape. Figure~\ref{fig:nvs} compares the methods on
a test image of Ryugu. All renderings reproduce the global appearance but differ in detail. PGSR reproduces the image
best, \pcgs\ renders the boulders with the correct shading but without cast shadows, and Asteroid-GS renders a nearly
uniform surface, because its geometry is smooth and its per-image appearance network has no embedding for an unseen
image. Over all test images, PGSR has the highest PSNR of the surface reconstruction methods, 39.3~dB on Itokawa and
35.2~dB on Ryugu, compared with 36.5 and 31.3~dB for \pcgs\ (Fig.~\ref{fig:nvs}c and Table~\ref{tab:images_full}), yet
its shapes are less accurate. Likewise, GOF reaches 38.0~dB on Itokawa with a mean distance of 2.40~m. \pcgs\ has no
appearance parameters with which to fit residual brightness variations, and its renderings of test views contain no
cast shadows, which leaves positive residuals along the shadowed limbs (Fig.~\ref{fig:nvs}d). Mip-Splatting illustrates
the point most clearly, with 42.6~dB on Itokawa and a mean distance of 2.48~m (Tables~\ref{tab:shape_full}
and~\ref{tab:images_full}). An appearance model with enough freedom can reproduce the images from an inaccurate surface,
and test-image metrics alone are therefore not suitable for ranking methods for shape modelling.

\begin{table*}[t]
\centering
\caption{Image metrics on the test images of Itokawa and Ryugu. Mean $\pm$ standard deviation over three splits.}
\label{tab:images_full}
\footnotesize
\setlength{\tabcolsep}{4.0pt}
\begin{tabular}{@{}lr@{\,$\pm$\,}lr@{\,$\pm$\,}lr@{\,$\pm$\,}lr@{\,$\pm$\,}l@{}}
\toprule
Method & \multicolumn{2}{c}{PSNR raw (dB)} & \multicolumn{2}{c}{PSNR (dB)} & \multicolumn{2}{c}{SSIM} & \multicolumn{2}{c}{LPIPS} \\
\midrule
\multicolumn{9}{@{}l}{\textit{Itokawa (AMICA, phase 10\textdegree{})}} \\
2DGS & 33.1 & 1.1 & 33.6 & 1.0 & 0.968 & 0.005 & 0.109 & 0.011 \\
GOF & 38.0 & 0.6 & 38.0 & 0.6 & 0.973 & 0.002 & 0.109 & 0.011 \\
PGSR & \underline{39.2} & 0.8 & \underline{39.3} & 0.8 & \underline{0.982} & 0.001 & \textbf{0.096} & 0.008 \\
Asteroid-GS & 16.7 & 0.8 & 29.6 & 0.9 & 0.926 & 0.010 & 0.144 & 0.009 \\
Asteroid-GS$^{\dagger}$ & 17.0 & 0.9 & 29.5 & 0.5 & 0.943 & 0.004 & 0.138 & 0.011 \\
Mip-Splatting$^{\S}$ & \multicolumn{2}{c}{\textbf{42.5}} & \multicolumn{2}{c}{\textbf{42.6}} & \multicolumn{2}{c}{\textbf{0.985}} & \multicolumn{2}{c}{0.100} \\
\textbf{PCGS (this work)} & 36.4 & 0.6 & 36.5 & 0.5 & 0.980 & 0.001 & \underline{0.098} & 0.008 \\
\addlinespace[3pt]
\multicolumn{9}{@{}l}{\textit{Ryugu (ONC-T, phase 19\textdegree{})}} \\
2DGS & \underline{34.8} & 1.8 & \underline{34.9} & 1.8 & \underline{0.948} & 0.005 & \underline{0.225} & 0.001 \\
GOF & 33.5 & 0.0 & 33.5 & 0.0 & 0.944 & 0.001 & \textbf{0.202} & 0.008 \\
PGSR & \textbf{35.1} & 0.2 & \textbf{35.2} & 0.2 & 0.947 & 0.003 & 0.227 & 0.002 \\
Asteroid-GS & 13.4 & 0.6 & 22.7 & 0.2 & 0.858 & 0.002 & 0.283 & 0.002 \\
Asteroid-GS$^{\dagger}$ & 13.5 & 0.6 & 22.7 & 0.0 & 0.868 & 0.004 & 0.282 & 0.002 \\
\textbf{PCGS (this work)} & 31.0 & 0.1 & 31.3 & 0.1 & \textbf{0.952} & 0.001 & \underline{0.225} & 0.002 \\
\bottomrule
\end{tabular}
\tablefoot{PSNR raw: without the per-image gain and offset. All other columns are computed after fitting one gain and offset per image over the body, for every method; methods with per-image appearance codes (Asteroid-GS, GOF) have no code for a test image and are rendered at a canonical exposure. $^{\dagger}$Asteroid-GS with the scene extent that its released code derives from the camera positions; the other rows use the object extent. $^{\S}$First split only. Bold: best; underlined: second best. Higher is better except for LPIPS.}
\end{table*}

%--------------------------------------------------------------------
\section{Discussion}
\label{sec:discussion}

\subsection{Accuracy in context}

Relative to the image resolution, the mean errors of \pcgs\ are at or below one pixel: 0.78 pixel on Itokawa, where the
pixel scale is 0.76~m, 0.28 pixel on Ryugu, where it is 2.2~m, and between 0.24 and 1.0 pixel on Bennu, whose images
range from 0.62 to 2.6~m per pixel. Mean errors below one pixel are possible because every surface element is seen in
many images, but they do not imply that features smaller than a pixel are resolved (Fig.~\ref{fig:recon}j). On
Itokawa, the error is close to one pixel and comparable to the mean facet size of the SPC reference itself, 0.60~m.

The large errors near the poles of Ryugu are most likely caused by the observation geometry. All five image sets were
acquired with the Sun within 8$^\circ$ of the equatorial plane of the body, at sub-solar latitudes between $-8^\circ$
and $+3^\circ$, and from near-equatorial viewpoints (Table~\ref{tab:data}). The polar regions are thus seen at grazing
emission and illuminated at grazing incidence in every image, and neither multi-view consistency nor shading
constrains them well. Since all methods show the same increase of the error towards the poles
(Fig.~\ref{fig:curves}f), the accuracy of \pcgs\ at high latitudes appears to be limited by the observation geometry
rather than by the method.

\subsection{Comparison with published results}
\label{sec:published}

\citet{Zhang2025AsteroidGS} reported mean distances of 0.705~m on Itokawa and 1.341~m on Ryugu for Asteroid-GS under
their own protocol. With the public code and its scene extent, but with our images and registration, we obtain 0.94~m on
Itokawa and 1.29~m on Ryugu. The Ryugu values agree; the difference on Itokawa cannot be traced because the image lists
and the evaluation code of that study are not public. \citet{ChenH2025} reported that 80.7\% and 95.5\% of the vertices
of their neural implicit model of Ryugu lie within 1 and 2~m of the SfM-MVS reference. Their models used 61 or 70 ONC-T
images, including a set at 0.6--0.7~m per pixel, and were explicitly supervised with SfM-MVS points. With the same
definition, 82.5\% and 97.1\% of the vertices of \pcgs\ lie within 1 and 2~m, obtained from 50 images at 2.2~m per pixel.
Since the protocols differ, these figures indicate comparable accuracy rather than a ranking.

\subsection{Role of the illumination}

The photometric model adds information where multi-view matching provides the least. Multi-view consistency fixes the
depth of textured surfaces, but in sparse image sets with narrow baselines many regions are smooth or seen from few
directions. In these regions, the shading under several Sun directions constrains the orientation of the surface, much
as brightness constrains the slopes between stereo control points in SPC \citep{Gaskell2008}. The largest gains are
thus expected on smooth terrain. Consistent with this expectation, the improvement over the best competing method is
larger on Itokawa, whose surface includes the smooth Muses-C region, than on Ryugu, with reductions in mean distance of
37\% and 15\%, respectively. However, the two bodies also differ in pixel scale and in their reference models, and the
Itokawa reference relies on the same photometric assumption as \pcgs. The difference between the two bodies is
therefore consistent with the expectation but does not demonstrate it.

The separation between albedo and shading also has limits. In the decomposition of Fig.~\ref{fig:overview}c, the
rendered albedo retains dark patches next to boulders, which indicates that part of the shading and of the short
shadows at the boulder scale is absorbed by the albedo instead of being explained by the geometry. The photometric
constraint thus appears to act mainly at scales of several pixels and above.

As noted for Bennu (Sect.~\ref{sec:bennu}), the comparison with the baselines does not isolate the photometric model.
Besides the rendering, \pcgs\ differs from the baselines in its surface extraction, its pixel convention, and its
night-side band, and each of these components matters (Sect.~\ref{sec:ablation}). The pixel-convention correction alone
reduced the mean distance on Ryugu from 1.35 to 0.76~m, more than the difference between \pcgs\ and PGSR, although the
alignment to the tie points removes the global part of this offset for the baselines (Sect.~\ref{sec:protocol}). The
ablation on Itokawa (Table~\ref{tab:development}) indicates that the shaded albedo improves on free view-dependent
colour when both are combined with multi-view consistency. Quantifying its contribution to the final model requires a
controlled comparison in which only the photometric model is replaced by free view-dependent colour.

\subsection{Limitations}

The Itokawa and Ryugu references are derived from images and may share systematic errors with image-based methods. The
Itokawa reference was built by stereophotoclinometry, which relies on the same photometric assumption as \pcgs, and the
Ryugu reference was built from images taken at 5--7~km \citep{ChenH2025}, which do not overlap with ours. The
references also have their own errors, which are largest where their source data were sparse, for instance near the
poles of Ryugu \citep{ChenH2025}. Differences between methods at high latitudes (Fig.~\ref{fig:curves}f) should
therefore be interpreted with caution. The Bennu reference removes the dependence on images for one body.

The scale of every reconstruction is fixed by the registration to the reference. With the SPICE scale alone, the volume
errors range from $-3.3$\% on Itokawa to $+1.0$\% on Ryugu at 40$^\circ$ (Table~\ref{tab:volume}), and a bulk density
derived from a first shape model would inherit this uncertainty. Regions that no image constrains, such as the poles of
Ryugu in our near-equatorial images and the deep neck of Itokawa, are closed by a smooth surface and should be flagged
as unmeasured in derived products.

The photometric model is a single lunar-Lambert function per image set, without a phase function or per-image exposure
terms, and albedo variations at the scales suppressed by the total-variation prior, such as the brighter Muses-C region
of Itokawa, are not represented explicitly. A photometric model with a phase function, fitted jointly across phase
angles \citep{Hapke2012}, could replace the switch between the two priors by a continuous description and would no
longer require $L$ to absorb the phase dependence of the brightness on Bennu. Several thresholds are set in metres
rather than in pixels (Table~\ref{tab:settings}); across our data they correspond to 0.2--5 pixels and would have to be
rescaled for images of other resolutions.

The phase threshold and the night-side band were chosen with the help of image sets that also serve for the evaluation,
and the comparison with the baselines is incomplete at 7$^\circ$ and 40$^\circ$ on Ryugu; the 7$^\circ$ image set is the
only one not used during development (Sect.~\ref{sec:images}). Finally, all image sets are hovering sequences that cover
all longitudes and have SPICE-based poses. Shorter arcs, as in a flyby, and less accurate poses remain to be tested.

\subsection{Implications for mission operations}

Training takes about 30~min on a single GPU, comparable to the other Gaussian splatting methods, and requires only the
images, their SPICE geometry, and a bundle adjustment of the camera poses. These properties make \pcgs\ a candidate for
first shape models early in a rendezvous. The scale can be taken from the trajectory, which in our data was accurate to
about 1\% in length and 3\% in volume, provided that unconstrained polar regions are flagged. Our results also suggest
how such first image sets could be acquired. Moderate phase angles may be favourable, as suggested by the most accurate
Ryugu model, obtained at 19$^\circ$. At higher phase angles, the night side must be treated explicitly, and views from
higher latitudes, which none of our image sets provides, would be needed to constrain the polar regions. Ongoing and
upcoming missions to small bodies, such as Hera \citep{Michel2022}, will acquire image sets of this kind.

%--------------------------------------------------------------------
\section{Conclusions}
\label{sec:conclusions}

We have presented photometrically constrained Gaussian splatting (\pcgs), which reconstructs watertight global shape
models of asteroids from a few dozen spacecraft images by incorporating the known solar illumination into Gaussian
splatting. Our main findings are as follows:
\begin{enumerate}
\item On Itokawa and Ryugu, the mean distance between the \pcgs\ models and the mission shape models is 0.59 and
0.62~m, 37\% and 15\% lower than that of the best of four Gaussian splatting methods. At the registered scale, the
volumes agree with the references to within 0.3\%.
\item Relative to the laser-altimeter model of Bennu, which is independent of the images, the mean distance is
$0.62\pm0.06$~m, whereas the other methods exceed 3.7~m. Most of this difference results from the treatment of the night
side.
\item Excluding a band beyond the terminator from the silhouette loss is essential at high phase angle and has little
effect on the mean distance at low phase angle. The phase-dependent prior improves the models of Itokawa, of Ryugu at
19$^\circ$, and of Bennu, but not those of Ryugu at 7$^\circ$ and 40$^\circ$.
\item The learned lunar-Lambert parameter of Ryugu decreases with increasing phase angle, from 0.78 at 7$^\circ$ to
0.37 at 40$^\circ$.
\item Fidelity to test images is not a reliable indicator of shape accuracy.
\end{enumerate}
Natural next steps are a controlled comparison with free view-dependent colour, a photometric model fitted jointly
across phase angles, tests on shorter image arcs and other bodies, and the application of the method to the first
images of upcoming missions.

\begin{acknowledgements}
We thank JAXA/ISAS and the Hayabusa and Hayabusa2 project teams, and NASA and the OSIRIS-REx team, for the images,
obtained from DARTS and the NASA Planetary Data System Small Bodies Node; R.~W.~Gaskell, S.~Watanabe and co-workers,
and the OSIRIS-REx Laser Altimeter team for the shape models; the NASA Navigation and Ancillary Information Facility
for the SPICE toolkit and kernels; and the authors of 2DGS, GOF, PGSR, and Asteroid-GS for releasing their code.
\end{acknowledgements}

\section*{Data availability}
The image lists, training and held-out splits, bundle-adjusted poses, masks, reconstructed meshes, and the code of \pcgs\
and of the evaluation will be made publicly available upon publication, with an archived copy on Zenodo. The AMICA,
ONC-T, and MapCam images are available from JAXA DARTS and the NASA Planetary Data System.

\bibliographystyle{aa}
\bibliography{refs}

\end{document}
